\section{The O2H Toolkit}\label{sec:o2h}
  Ever since their first introduction by Unruh~\cite{EC:Unruh14}, the One-Way
    to Hiding (O2H) lemmas
    have served as indispensible tools for cryptographers seeking to prove
    security in the quantum random oracle model.
    Note the plural: Over the years, the original O2H lemma, which only
    introduced a very limited capability of reprogramming QROs, has
    been repeatedly generalized and expanded, often to different
    and seemingly-incomparable settings. Thus, there is not one
    ``multi-tool'' O2H Lemma, but several
    variants that together make up a ``toolkit''---and, as with any good
    toolkit, different tools are best fit for different jobs!

  In the following sections, we will introduce the different variants in
    increasing generality and complexity, starting with the very simplest
    variant, and showing along the way how to apply each to a security
    reduction. An overview of the O2H lemmas that will be presented is provided in
    \autoref{tab:o2h}.\footnote{
      There are other historical variants of the O2H lemma that are not covered here,
      due later variants being strict generalisations. These include (\ldots)
    }

  \begin{table}[h]
    \centering
    %
    \caption{
      \label{tab:o2h}
      An overview of the various One-way to Hiding (O2H) lemmas. 
      Here, Q $=$ Quantum-accessible oracles, M $=$ Multi-point reprogramming, A $=$ Adaptive
      reprogramming, B $=$ fully (as opposed to rewindable) Black-box access to
      the distinguisher, C $=$ Computationally (as opposed to
      statistically) hidden reprogramming, L $=$ allows Leaking information about
      reprogrammed points after reprogramming, and G $=$ Guess, meaning the
      extractor does not have access to an indicator function for the target
      value(s). $q$ bounds the total number of
      oracle calls, and $r$ bounds the number of reprogrammed points. 
      Variants in paranthesis have been superseded by later results, and are
      included mainly for historical reasons.
      Most bounds presented here are simplified and/or approximate, removing
      sums and ignoring constant factors and information-theoretic terms; for
      the exact statement, see the respective Lemma. 
      Finally, to make for easier comparison, the
      ``Set $\secpar$ to'' column gives a numerical example for how much one
      would have to increase the security parameter if one wants to bound the
      distinguishing advantage by $2^{-128}$, assuming 1) that the extractors
      advantage is bounded by $\varepsilon \leq 2^{-\secpar}$, 2) $q \leq 2^{80}$ queries
      to the random oracle (a common estimate for the total number of SHA3
      evaluations over the course of its lifetime), and 3) $r \leq 2^{10}$
      reprogrammed points (following the example of Bellare et
      al.~\cite{AC:BRTZ24}, who bounded the number of user corruptions
      similarly); the closer to the optimal value $\secpar = 128$ the better.
    }
    %
    \rowcolors{2}{}{gray!10}
    \begin{tabular}{ccccccccccc}
      \toprule
      Ref. & O2H Variant  & \rotatebox{90}{Guessing} 
                          & \rotatebox{90}{Quantum} 
                          & \rotatebox{90}{Multi-point} 
                          & \rotatebox{90}{Adaptive} 
                          & \rotatebox{90}{Computational} 
                          & \rotatebox{90}{Leaking} 
                          & \rotatebox{90}{Black-box} 
                          & Bound (simplified) & Set $\secpar$ to \\
      \midrule
      \cite{CCS:BelRog93} &
        Classical O2H (guess) 
        & \Tyes & \Tno & \Tyes & \Tyes & \Tyes & \Tyes & \Tyes &
        $\leq q \varepsilon$ & 208 \\
      \cite{CCS:BelRog93} &
        Classical O2H (search) 
        & \Tno & \Tno & \Tyes & \Tyes & \Tyes & \Tyes & \Tyes & 
        $\leq \varepsilon$ & 128 \\
      \midrule
      \cite{EC:Unruh14} &
        Original O2H 
        & \Tyes & \Tyes & \Tno & \Tno & \Tyes & \Tyes & \Tyes &
        $\lesssim q \sqrt{\varepsilon}$ & 416 \\
      \cite{C:Unruh14} &
        ``Adaptive'' O2H 
        & \Tyes & \Tyes & \Tno & \Tyes & \Tyes & \Tyes & \Tyes &
        $\lesssim q \sqrt{\varepsilon}$ & 416 \\
      \midrule
      \cite{C:AmbHamUnr19} &
        ``Semi-Classical'' O2H (guess) 
        & \Tyes & \Tyes & \Tyes & \Tno & \Tyes & \Tyes & \Tyes &
        $\lesssim q \sqrt{\varepsilon}$ & 416 \\
      \cite{C:AmbHamUnr19} &
        ``Semi-Classical'' O2H (search) 
        & \Tno & \Tyes & \Tyes & \Tno & \Tyes & \Tyes & \Tyes &
        $\lesssim \sqrt{q \varepsilon}$ & 336 \\
      \midrule
      \cite{AC:GHHM21} &
        Resampling Lemma 
        & \Tyes & \Tyes & \Tyes & \Tyes & \Tno & \Tyes & \Tyes &
        $\lesssim r \sqrt{q \varepsilon}$ & 356 \\
      %Fixed-Perm.\ RL & \Tyes & \Tyes & \Tyes & \Tno & \Tyes & \Tyes &
      %  $\advantage{\dist}{e}[\Dis{}] \lesssim \qRP \sqrt{ \qRO \cdot \advantage{\guess}{}[\Ext{}] }$ & 
      %  $\nicefrac{1}{\varepsilon} \approx \qRP^2 \qRO$ & \autoref{lemma:resampling-fp} & \cite{EPRINT:Jaeger24b} \\
      %\midrule
      \cite{PKC:PanZen24} &
        Comp.\ Ad.\ Reprog.\ Framew. 
        & \Tyes & \Tyes & \Tyes & \Tyes & \Tyes & \Tno & \Tyes &
        $\lesssim r q \sqrt{\varepsilon}$ & 436 \\
      \cite{EPRINT:Jaeger24b} &
        ``Fixed-Permutation'' O2H
        & \Tyes & \Tyes & \Tyes & \Tyes & \Tyes? & \Tyes? & \Tyes &
        $\lesssim \sqrt{q\varepsilon}$? & 336? \\
      \midrule
      \cite{TCC:BHHHP19} &
        ``Double-sided'' O2H 
        & \Tno & \Tyes & \Tno & \Tno & \Tyes & \Tyes & \Tyes &
        $\lesssim \sqrt{\varepsilon}$ & 256 \\
      \midrule
      \cite{EC:KSSSS20} &
        Meas.-Rewind-Meas. 
        & \Tno & \Tyes & \Tyes & \Tno & \Tyes & \Tyes & \Tno &
        $\lesssim q \varepsilon$ & 208 \\
      \cite{AC:GeLiaXue24} &
        Meas.-Rewind-Ext. 
        & \Tno & \Tyes & \Tyes & \Tno & \Tyes & \Tyes & \Tno &
        $\lesssim \sqrt{q} \varepsilon$ & 168 \\
      \bottomrule
    \end{tabular}
  \end{table}

  \subsection{Guess versus Find}
    Let us show that the extractor having access to both $\RO$ and $\GO$ is
      equivalent to the extractor having access to a quantum-accessible
      indicator function.
    
    \begin{lemma}[Double-sided oracle access $\Leftrightarrow$ Indicator oracle access]
      Let $\careset \subseteq \Xspace$, let $\RO, \GO: \Xspace \to \Yspace$
      be classical random oracles satisfying $\forall x \not\in \careset: \RO(x) =
      \GO(x)$, and let $\indfunc$ be an indicator oracle for $\careset$. 
      A quantum oracle algorithm with oracle access $\indfunc$ can
      faithfully simulate $\RO$ and $\GO$, and a quantum oracle algorithm with
      oracle access to $\RO$ and $\GO$ can simulate $\indfunc$, up to collisions in
      $\careset$.
    \end{lemma}

    \begin{proof}[sketch]
      First direction: The algorithm maintains two independent simulations of
      random oracles using the compressed random oracle technique; call them
      $\RO_1$ and $\RO_2$. Upon receiving a query to $\RO$, the algorithm
      responds using $\RO_1$. Upon receiving a query to $\GO$, the algorithm
      computes the output of both $\RO_1$ and $\RO_2$, storing each in separate output
      registrers, then uses $\indfunc$ to coherently implement a control for a
      CSWAP operation that swaps the first and second output register. The simulated
      $\GO$ then behaves like $\RO_1$ on inputs $x \in \careset$, and like
      $\RO_2$ on inputs $x \not\in \careset$.

      Second direction: Recall that the indicator oracle operates on an $n$-length
      input register and a single-qubit output register, which is flipped if the
      input $x$ is in $\careset$. To simulate $\indfunc$,
      the algorithm pairs the input register with an $m$-length output register
      initilazied to $\ket{0^m}$, then first calls $\RO$, followed by $\GO$. The
      resulting state will have basis states $\ket{x, \RO(x) \xor \GO(x)}$.
      Using CNOTs to implement an OR-function, the output bit is then flipped
      iff the state of the output is not equal to $\ket{0^m}$. Thus, we get an
      oracle that indicates whether $\RO(x)$ is unequal to $\GO(x)$, which is
      identical to the indicator function for $\careset$ up to collisions in
      $\careset$ (for which the simulation will return a false negative).
      Before returning the result, the state of the output register is
      uncomputed via one more call each to $\RO$ and $\GO$, removing 
      unwanted entanglement.
      \hhnote{
        To-do: Quantify both the loss due to collisions and due to any residue
        entanglement.
      }
    \end{proof}

  \subsection{Non-Adaptive State of the Art}
    In the lemmas below, $\RO$, $\GO$, $\careset$ (or $\target$), and $z$ are random
      with arbitrary joint distributions.

    \begin{lemma}[Classical O2H~\cite{CCS:BelRog93}]
      Let $\careset \subseteq \Xspace$, let $\RO, \GO: \Xspace \to \Yspace$
      be classical random oracles satisfying $\forall x \not\in \careset: \RO(x) =
      \GO(x)$, let $\indfunc$ be an indicator oracle for $\careset$, and let $z$ be arbitrary. Then 
      there are SFBB extractors $\findext, \guessext$ such that, for all distinguishers
      $\distinguisher$ making at most $\qRO$ calls to its random oracle,
      \[
        \abs{
          \prob{1 \sample \distinguisher^{\RO}(z)} - \prob{1 \sample \distinguisher^{\GO}(z)}
        } \leq \probsample{x \sample \findext^{\indfunc}(z)}{x \in \careset}\, ,
      \]
      and
      \[
        \abs{
          \prob{1 \sample \distinguisher^{\RO}(z)} - \prob{1 \sample \distinguisher^{\GO}(z)}
        } \leq \qRO \cdot \probsample{x \sample \guessext(z)}{x \in \careset}\, .
      \]
      The extractors run in time that of $\distinguisher$ plus the
      overhead of simulating at most $\qRO$ oracle calls, and
      $\findext$ additionally makes at most $\qRO$ calls to its
      indicator oracle.
    \end{lemma}

    \begin{lemma}[Original O2H Lemma~\cite{EC:Unruh14}]
      Let $\RO, \GO: \Xspace \to \Yspace$
      be random oracles satisfying $\forall x \neq \target: \RO(x) =
      \GO(x)$ (and independent otherwise), and let $z$ be a value that depends arbitrarily on $\RO, \GO$,
      and $\target$. Then there is an SFBB extractor $\extractor$ such that, for all distinguishers
      $\distinguisher$ making at most $\qRO$ calls to its random oracle,
      \[
        \abs{
          \prob{1 \sample \distinguisher^{\RO}(z)} - \prob{1 \sample \distinguisher^{\GO}(z)}
        } \leq 2\qRO \sqrt{\prob{\target \sample \extractor(z)}}\, ,
      \]
      where the extractor runs in time that of $\distinguisher$, plus the
      overhead of simulating at most $\qRO$ oracle responses.
    \end{lemma}

    \begin{lemma}[Multi-point O2H~\cite{C:AmbHamUnr19}]
      Let $\careset \subseteq \Xspace$, let $\RO, \GO: \Xspace \to \Yspace$
      be quantum random oracles satisfying $\forall x \not\in \careset: \RO(x) =
      \GO(x)$, let $\indfunc$ be a quantum-accessible indicator oracle for $\careset$, and let $z$ be arbitrary. Then 
      there are SFBB extractors $\findext, \guessext$ such that, for all distinguishers
      $\distinguisher$ making at most $\qRO$ calls to its random oracle,
      $\forall e \in \{1, \half\}$,
      \[
        \abs{
          \prob{1 \sample \distinguisher^{\RO}(z)}^e - \prob{1 \sample \distinguisher^{\GO}(z)}^e
        } \leq 2 \sqrt{
          (\qRO + 1) \cdot \probsample{x \sample \findext^{\indfunc}(z)}{x \in \careset}
        }\, ,
      \]
      and
      \[
        \abs{
          \prob{1 \sample \distinguisher^{\RO}(z)}^e - \prob{1 \sample \distinguisher^{\GO}(z)}^e
        } \leq 2 \qRO \sqrt{
          \probsample{x \sample \guessext(z)}{x \in \careset}
        }\, .
      \]
      The extractors run in time that of $\distinguisher$ plus the
      overhead of simulating at most $\qRO$ oracle responses, and
      $\findext$ additionally makes at most $\qRO$ calls to its
      indicator oracle.
    \end{lemma}

    \begin{lemma}[Single-point O2H, find~\cite{TCC:BHHHP19}]
      Let $\target \in \Xspace$, let $\RO, \GO: \Xspace \to \Yspace$
      be quantum random oracles satisfying $\forall x \neq \target: \RO(x) =
      \GO(x)$, let $\indfunc$ be a quantum-accessible indicator oracle for $\target$, and let $z$ be arbitrary. Then 
      there is an SFBB extractor $\extractor$ such that, for all distinguishers
      $\distinguisher$ making at most $\qRO$ calls to its random oracle, 
      $\forall e \in \{1, \half\}$,
      \[
        \abs{
          \prob{1 \sample \distinguisher^{\RO}(z)}^e - \prob{1 \sample \distinguisher^{\GO}(z)}^e
        } \leq 2 \sqrt{
          \prob{\target \sample \extractor^{\indfunc}(z)}
        }\, ,
      \]
      where $\extractor$ runs in time that of $\distinguisher$ plus the
      overhead of simulating at most $2\qRO$ oracle responses, and makes at most $\qRO$ calls to its
      indicator oracle.
    \end{lemma}

    \begin{lemma}[Multi-point O2H with rewinding, find~\cite{AC:GeLiaXue24}]
      Let $\careset \subseteq \Xspace$, let $\RO, \GO: \Xspace \to \Yspace$
      be quantum random oracles satisfying $\forall x \not\in \careset: \RO(x) =
      \GO(x)$, let $\indfunc$ be a quantum-accessible indicator oracle for
      $\careset$, and let $z$ be arbitrary. Then 
      there is an RBB extractor $\extractor$ such that, for all distinguishers
      $\distinguisher$ making at most $\qRO$ calls to its random oracle, 
      \[
        \abs{
          \prob{1 \sample \distinguisher^{\RO}(z)} - \prob{1 \sample \distinguisher^{\GO}(z)}
        } \leq 4 \sqrt{\qRO} \cdot \probsample{x \sample \extractor^{\indfunc}(z)}{x \in \careset}\, ,
      \]
      where $\extractor$ runs in time thrice that of $\distinguisher$ (rewinding
      twice), plus the overhead of simulating at most $6\qRO$ oracle responses,
      and makes at most $3\qRO$ calls to its indicator oracle.
    \end{lemma}

    \hhnote{
      MRM required access to both $\RO$ and $\GO$; exchanged for indicator
      function, making MRE strictly more general. Also double-check that it is
      rewinding twice.
    }

    \hhnote{Attempt at putting them all in one lemma}
    \begin{lemma}[Non-adaptive O2H, state of the art]
      Let $\RO, \GO: \Xspace \to \Yspace$
      be quantum random oracles, let $\careset \subseteq \Xspace$ be such that
      $\forall x \not\in \careset: \RO(x) = \GO(x)$, let $\indfunc$ be a
      quantum-accessible indicator oracle for $\careset$, and let $z$ be a value
      that depends arbitrarily on $\RO, \GO$ and $\careset$.
      Let $\distinguisher$ be a distinguisher that makes at most $\qRO$ calls to
      its random oracle. Then,
      \begin{enumerate}
        \item \textbf{(Classical guess)} There is an SFBB extractor $\extractor$ such that, for all $\distinguisher$ that only
          make classical calls to the random oracle, 
          \[
            \abs{
              \prob{1 \sample \distinguisher^{\RO}(z)} - \prob{1 \sample \distinguisher^{\GO}(z)}
            } \leq \qRO \cdot \probsample{x \sample \extractor(z)}{x \in \careset}\, ,
          \]
          where $\extractor$ runs in time that of $\distinguisher$ plus the
          overhead of simulating at most $\qRO$ oracle responses.
        \item \textbf{(Classical search)} There is an SFBB extractor $\extractor$ such that, for all $\distinguisher$ that only
          make classical calls to the random oracle, 
          \[
            \abs{
              \prob{1 \sample \distinguisher^{\RO}(z)} - \prob{1 \sample \distinguisher^{\GO}(z)}
            } \leq \probsample{x \sample \extractor^{\indfunc}(z)}{x \in \careset}\, ,
          \]
          where $\extractor$ runs in time that of $\distinguisher$ plus the
          overhead of simulating at most $\qRO$ oracle responses,
          and makes at most $\qRO$ calls to its indicator oracle.
        \item \textbf{(Multi-point guess)} There is an SFBB extractor $\extractor$ such
          that, for all $\distinguisher$ and $e \in \{1, \half\}$,
          \[
            \abs{
              \prob{1 \sample \distinguisher^{\RO}(z)}^e - \prob{1 \sample \distinguisher^{\GO}(z)}^e
            } \leq 2 \qRO \sqrt{
              \probsample{x \sample \extractor(z)}{x \in \careset}
            }\, ,
          \]
          where $\extractor$ runs in time that of $\distinguisher$ plus the
          overhead of simulating at most $\qRO$ oracle responses.
        \item \textbf{(Multi-point search)} There is an SFBB extractor $\extractor$ such
          that, for all $\distinguisher$ and $e \in \{1, \half\}$,
          \[
            \abs{
              \prob{1 \sample \distinguisher^{\RO}(z)}^e - \prob{1 \sample \distinguisher^{\GO}(z)}^e
            } \leq 2 \sqrt{
              (\qRO + 1) \cdot \probsample{x \sample \extractor^{\indfunc}(z)}{x \in \careset}
            }\, ,
          \]
          where $\extractor$ runs in time that of $\distinguisher$ plus the
          overhead of simulating at most $\qRO$ oracle responses,
          and makes at most $\qRO$ calls to its indicator oracle.
        \item \textbf{(Single-point search)} For $\careset = \{\target\}$, there is an SFBB extractor $\extractor$ such
          that, for all $\distinguisher$ and $e \in \{1, \half\}$,
          \[
            \abs{
              \prob{1 \sample \distinguisher^{\RO}(z)}^e - \prob{1 \sample \distinguisher^{\GO}(z)}^e
            } \leq 2 \sqrt{
              \prob{\target \sample \extractor^{\indfunc}(z)}
            }\, ,
          \]
          where $\extractor$ runs in time that of $\distinguisher$ plus the
          overhead of simulating at most $2\qRO$ oracle responses,
          and makes at most $\qRO$ calls to its indicator oracle.
        \item \textbf{(Multi-point search with rewinding)} There is an RBB
          extractor $\extractor$ such that, for all $\distinguisher$,
          \[
            \abs{
              \prob{1 \sample \distinguisher^{\RO}(z)} - \prob{1 \sample \distinguisher^{\GO}(z)}
            } \leq 4 \sqrt{\qRO} \cdot \probsample{x \sample \extractor^{\indfunc}(z)}{x \in \careset}\, ,
          \]
          where $\extractor$ runs in time at most thrice that of $\distinguisher$ (rewinding
          twice), plus the overhead of simulating at most $6\qRO$ oracle
          responses, and makes at most $3\qRO$ calls to its indicator oracle.
      \end{enumerate}
    \end{lemma}

  \subsection{Making the Lemmas Adaptive}

  % Subsection: Classical O2H
    %\section{Classical O2H and Bellare--Rogaway Encryption}
  \hhnote{
    In which I define our motivating scheme, define and prove classical
    variants of the O2H lemma, and show the classical (multi-chal.) security of BRE.
  }



  %% Extract the following into separate files and sections
      % as you get to them
  %\subsection{Single-Challenge Implication from the O2H Lemma}
  %  \hhnote{
  %    In which I introduce the original version of O2H, and use it to prove
  %    our first result: Single-challenge $\owcpa$ KEM yields $\indcpa$ PKE.
  %  }

  %\subsection{Multi-Challenge Implication from Semi-Classical O2H}
  %  \hhnote{
  %    In which I introduce the Semi-Classical O2H, and use it to prove
  %    that if we start from $\owpca$, we get something tighter than applying 
  %    the generic hybrid reduction to the previous theorem.
  %  }

  %\subsection{Tightening the Bound with the Resampling Lemma}
  %  \hhnote{
  %    In which I introduce the Resampling Lemma, and apply it to the first
  %    game hop to make it ever so slightly tighter.
  %  }

  %\subsection{Tighter Still from qPCA}
  %    \hhnote{
  %      In which I introduce qPCA, and show that if we have this, then we can
  %      exploit the modularity of the Semi-Classical O2H theorems to get a
  %      tighter result.
  %    }

  %  \paragraph{Moving between stronger notions.}
  %  \hhnote{
  %    In which I note that (with proof sketches) that what we have so far
  %    already suffices to move between CCA-notions, and even between notions
  %    with user corruptions, as long as opening challenges is disallowed,
  %    (since we can still assume that all the reprogrammed points were defined
  %    at the start of the game).
  %  }

  %\subsection{Adding Challenge Opening with Adaptive Reprogramming O2H}
  %  \hhnote{
  %    In which I introduce the notion of multi-bit, multi-challenge with sender
  %    openings, and argue that now (finally), the above techniques are
  %    insufficient. Introduce the Pan--Zeng adaptive reprogramming framework,
  %    and \emph{either} use it directly to show security, and then show how
  %    Joseph re-derived it from permutations, \emph{or} show how Joseph is able
  %    to get adaptivity from Semi-Classical O2H directly, and use this to prove
  %    the result (whichever one gives a simpler and more intuitive proof).
  %  }
